\usepackage{amsmath,amssymb,amsthm}
% \usepackage{amsthm}
% \usepackage[margin=2cm]{geometry}
\usepackage{dsfont}
\usepackage{tikz}
\usetikzlibrary{cd}
\usetikzlibrary{babel,calc}
\usetikzlibrary{positioning,arrows.meta,fit,backgrounds}
% \usepackage{tikz-cd}
% \usepackage{titlesec}
\usepackage{enumitem}
\usepackage{hyperref}
\usepackage{cleveref}
% \usepackage{parskip}
\usepackage{ebproof}
% \usepackage{quiver}
\usepackage{stackengine,scalerel}
% \usepackage{xparse}
% \usepackage{physics}
% \usepackage{physics2}
% \usephysicsmodule{ab,ab.braket}
% \usepackage{tikzit}
% \input{smc.tikzstyles}
\usepackage{xcolor}
\usepackage{graphicx}
\usepackage{array}    % m{} columns + \arraybackslash (combined car-model table)
\usepackage{booktabs} % \toprule/\midrule/\bottomrule

\usepackage{mathtools}
\usepackage{stmaryrd}
\usepackage{subcaption}

\newtheorem*{fact}{Fact}

\crefname{prop}{prop.}{props.}
\Crefname{prop}{Proposition}{Propositions} % Capitalized version for start of sentences

\crefname{remark}{remark}{remarks}
\Crefname{remark}{Remark}{Remarks} % Capitalized version for start of sentences

\crefname{appendix}{appendix}{appendices}
\Crefname{appendix}{Appendix}{Appendices} % Capitalized version for start of sentences

% \crefname{property}{property}{properties}
% \Crefname{property}{Property}{Properties} % Capitalized version for start of sentences

\crefformat{property}{#2#1#3}
\Crefformat{property}{#2#1#3}

\crefname{def}{definition}{definitions.}
\Crefname{def}{Definition}{Definitions} % Capitalized version for start of sentences

\crefformat{section}{#2\S#1#3}
\Crefformat{section}{#2\S#1#3}

\usepackage{listings}

\usepackage{tcolorbox}
\tcbuselibrary{listings} % Import listings support

\definecolor{codegreen}{rgb}{0,0.6,0}
\definecolor{codegray}{rgb}{0.5,0.5,0.5}
\definecolor{codepurple}{rgb}{0.58,0,0.82}
\definecolor{backcolour}{rgb}{0.95,0.95,0.92}

\lstdefinestyle{mystyle}{
    backgroundcolor=\color{backcolour},
    commentstyle=\color{codegreen}\textit,
    keywordstyle=\color{blue},
    numberstyle=\tiny\color{codegray},
    stringstyle=\color{codepurple},
    basicstyle=\ttfamily\footnotesize,
    breakatwhitespace=false,
    breaklines=false,
    captionpos=b,
    keepspaces=true,
    numbers=left,
    numbersep=5pt,
    showspaces=false,
    showstringspaces=false,
    showtabs=false,
    tabsize=2
}

\lstset{style=mystyle}

\tcbset{
    mystyle/.style={
        listing only,
        breakable, % Allows box to span pages
        colback=CodeBack,
        colframe=CodeFrame,
        boxrule=0.5pt,
        sharp corners,
        listing options={
            backgroundcolor=\color{backcolour},   
            commentstyle=\color{codegreen}\textit,
            keywordstyle=\color{blue},
            numberstyle=\tiny\color{codegray},
            stringstyle=\color{codepurple},
            basicstyle=\ttfamily\footnotesize,
            breakatwhitespace=false,         
            breaklines=true,                 
            captionpos=b,                    
            keepspaces=true,                 
            numbers=left,                    
            numbersep=5pt,                  
            showspaces=false,                
            showstringspaces=false,
            showtabs=false,                  
            tabsize=2,
        }
    }
}

\newtcblisting{codelisting}[3][]{
    listing only,
    % breakable, % Allows box to span pages
    colback=backcolour,
    % colframe=CodeFrame,
    boxrule=0pt,
    boxsep=0pt,
    top=1mm,
    bottom=1mm,
    sharp corners,
    notitle,
    listing options={
        language=#2,
        morekeywords= {exo, endo},
        commentstyle=\color{codegreen}\textit,
        keywordstyle=\color{blue},
        numberstyle=\tiny\color{codegray},
        stringstyle=\color{codepurple},
        basicstyle=\ttfamily\tiny,
        breakatwhitespace=false,         
        breaklines=true,                 
        captionpos=b,                    
        keepspaces=true,                 
        numbers=left,                    
        numbersep=5pt,                  
        showspaces=false,                
        showstringspaces=false,
        showtabs=false,                  
        tabsize=2,
        aboveskip=0pt,  % Removes the space above the code
        belowskip=0pt,  % Removes the space below the code
        mathescape=#3,
        #1               % optional extra listings options (e.g. escapeinside for \cnew)
    }
}

% Highlight code that is new relative to the previous rung's program.
% Use inside a codelisting that passes escapeinside, e.g.
%   \begin{codelisting}[escapeinside={(*}{*)}]{haskell}{true} ... (*\cnew{new}*) ...
\definecolor{codenew}{RGB}{255,228,130}
\newcommand{\cnew}[1]{{\setlength{\fboxsep}{0.6pt}\colorbox{codenew}{\ttfamily #1}}}
% Coloured (not highlighted) font, for marking new syntax in the grammar figure.
\newcommand{\cnewm}[1]{\textcolor{blue!65!black}{#1}}

% Theo's Macros
\newcommand{\naturals}{\mathbb{N}}
\newcommand{\integers}{\mathbb{Z}}
\newcommand{\booleans}{\mathbb{B}}
\newcommand{\indicator}{\mathds{1}}
\newcommand{\va}{\mathbf{a}}
\newcommand{\vb}{\mathbf{b}}
\newcommand{\vx}{\mathbf{x}}
\newcommand{\vy}{\mathbf{y}}
\newcommand{\vw}{\mathbf{w}}
\newcommand{\vz}{\mathbf{z}}
\newcommand{\vsigma}{\mathbf{\sigma}}
\newcommand{\reals}{\mathbb{R}}
\newcommand{\complex}{\mathbb{C}}
\newcommand{\field}{\mathbb{F}}
\newcommand{\sign}{\mathrm{sign}}
\DeclareMathOperator*{\argmin}{arg\,min}
\newcommand{\true}{\mathsf{True}}
\newcommand{\false}{\mathsf{False}}
\newcommand{\brackets}[1]{\{#1\}}

% \newcommand{\bicatb}{\mathcal{B}}

\newcommand{\catc}{\mathbf{C}}
\newcommand{\catcop}{\catc^{op}}
\newcommand{\catd}{\mathbf{D}}
\newcommand{\cate}{\mathcal{E}}
\newcommand{\eqdef}{\triangleq}
\newcommand{\id}{\mathbf{id}}
% \newcommand{\idbicat}{\mathbf{Id}}
\newcommand{\Id}{\mathbf{Id}}
\newcommand{\righttop}[1]{\xrightarrow{#1}}
\newcommand{\lefttop}[1]{\xleftarrow{#1}}
\newcommand{\Ob}{\mathbf{Ob} \ }
\newcommand{\Hom}{\mathrm{Hom}}
% \newcommand{\yo}{\textmc{よ}}
\newcommand{\yo}{\textbf{y}}

\newcommand{\terminalmorph}{\langle\rangle}
\newcommand{\prodmorph}[2]{\langle#1, #2\rangle}

\newcommand{\ul}[1]{\underline{#1}}
\newcommand{\ulyo}{\ul{\yo}}
\newcommand{\yozero}{\yo(0)}
\newcommand{\yoone}{\yo(1)}
\newcommand{\yon}{\yo(n)}
\newcommand{\yom}{\yo({m})}
\newcommand{\yop}{\yo(p)}
% \newcommand{\op}{\mathbf{op}}
\newcommand{\Ops}{\mathcal{O}}
% \newcommand{\typing}[3]{#1 \ | \ #2 \vdash #3}
% \newcommand{\typingfilled}[1]{\Gamma \ | \ \Delta \vdash #1}
% \newcommand\eye{\ensurestackMath{\stackinset{c}{}{c}{-.33pt}%
%   {\bullet}{\bigcirc}}}
% \newcommand\smallereye{\scalerel*{\eye}{x}}
% \newcommand{\eyemid}{ \ \eye \ }

% Vectors
\newcommand{\veca}{\Vec{a}}
\newcommand{\vecb}{\Vec{b}}
\newcommand{\vecc}{\Vec{c}}
\newcommand{\vecd}{\Vec{d}}
\newcommand{\vecx}{\Vec{x}}
\newcommand{\vecy}{\Vec{y}}
\newcommand{\vecz}{\Vec{z}}

\newcommand{\alphaequiv}{\simeq_\alpha}

\newcommand{\LHS}{\mathsf{LHS}}
\newcommand{\RHS}{\mathsf{RHS}}

% Chapter 2 algebraic theories
\newcommand{\Mon}{\mathsf{Mon}}

\newcommand{\Tree}{\mathsf{Tree}}

\newcommand{\Kl}[1]{\mathsf{Kl}({#1})}

\newcommand{\Name}{\mathcal{N}}
\newcommand{\Var}{\mathsf{Var}}
\newcommand{\IPVar}{\mathsf{IPVar}}
\newcommand{\GroundValue}{\mathsf{GroundValue}}
\newcommand{\Value}{\mathsf{Value}}
\newcommand{\Expr}{\mathsf{Expr}}

\newcommand{\bind}{\mathbin{>\!\!\!>\mkern-6.7mu=}}
\newcommand{\fin}{\mathsf{fin}}
\newcommand{\MVal}{\mathsf{MVal}}
\newcommand{\Intervention}{\mathsf{Intv}}
\newcommand{\MultiWorld}{\mathsf{MultiWorldCounterfactual}}
\newcommand{\TwinWorld}{\mathsf{TwinWorldCounterfactual}}
\newcommand{\supp}{\mathsf{supp}}

% ------------- Formal Def ----------------
\newcommand{\IdealisedChiRho}{\mathsf{ChiRho}}
\newcommand{\Prob}{\mathsf{Prob}}
\newcommand{\Dist}{\mathsf{Dist}}
\newcommand{\Paths}{\mathsf{Paths}}
\newcommand{\UPaths}{\mathsf{UPaths}}
\newcommand{\posreals}{\reals_{\geq 0}}

% Syntax
% \newcommand{\opif}[3]{\mathbf{if} \ #1 \ #2 \ #3}
% \newcommand{\lift}[1]{\overline{#1}}
\newcommand{\opprod}{\mathsf{prod}}
\newcommand{\opscore}{\mathsf{score}}
\newcommand{\opflip}{\mathsf{flip}}
\newcommand{\dbernoulli}{\mathbf{bernoulli}}
\newcommand{\dnormal}{\mathbf{normal}}
\newcommand{\constr}{\mathbf{constr}}
\newcommand{\intvnone}{\mathsf{None}}
\newcommand{\intvsome}{\mathsf{Some}}
\newcommand{\opskip}{\mathsf{skip}}
\newcommand{\opsample}{\mathsf{sample}}
\newcommand{\opintervene}{\mathsf{intervene}}
\newcommand{\opbroadcast}{\mathsf{broadcast}}
\newcommand{\optrack}{\mathsf{track}}
\newcommand{\opdistribute}{\mathsf{distribute}}
\newcommand{\opif}{\mathbf{if}}
\newcommand{\ifelse}[3]{\opif \ #1 \ \mathbf{then} \ #2 \ \mathbf{else} \ #3}
\newcommand{\opthunk}{\mathsf{thunk}}
\newcommand{\opforce}{\mathsf{force}}
\newcommand{\returnSymb}{\mathbf{return}}
\newcommand{\return}[1]{\returnSymb \ #1}
\newcommand{\returnParam}[2]{\returnSymb_{#1} \ #2}
\newcommand{\app}{\odot}
\newcommand{\pureval}{\mathbf{pureVal}}
\newcommand{\mem}{\mathsf{mem}}
\newcommand{\letbind}[3]{\mathbf{let} \ #1 = #2 \ \mathbf{in} \ #3}
\newcommand{\opdo}{\mathsf{do}}
\newcommand{\rhide}{\mathsf{hide}}
\newcommand{\rwk}{\mathsf{wk}}
\newcommand{\rcontr}{\mathsf{contr}}
\newcommand{\rexch}{\mathsf{exch}}
\newcommand{\rrename}{\mathsf{rename}}
\newcommand{\Renamings}{\mathsf{Renamings}}
\newcommand{\Deltahat}{\widehat{\Delta}}
\newcommand{\projf}{\pi_{\mathsf{f}}}
\newcommand{\projcf}{\pi_{\mathsf{cf}}}
\newcommand{\coerce}[3]{\mathsf{coerce}_{#2}^{#1}[#3]}
\newcommand{\rv}{\underline{r}}
\newcommand{\restr}{\rvert}
\newcommand{\subst}[2]{[#1/#2]}
\newcommand{\ket}[1]{| #1 \rangle}

% Types
\newcommand{\tnat}{\mathsf{Nat}}
\newcommand{\tbool}{\mathsf{Bool}}
\newcommand{\tint}{\mathsf{Int}}
\newcommand{\tposreal}{\mathsf{PosReal}}
\newcommand{\treal}{\mathsf{Real}}
\newcommand{\tunit}{\mathsf{Unit}}
\newcommand{\tprod}{\mathsf{Prod}}
\newcommand{\tP}[2]{P \ #1 \ #2}
\newcommand{\unitgrade}{[]}
\newcommand{\rarrow}{\Rightarrow}

% Type system
\newcommand{\dtyping}[3]{#1 \vdash^d #2: #3}
\newcommand{\ptyping}[4]{#1 \mid #2 \vdash^p #3: #4}
\newcommand{\rtyping}[4]{#1 \vdash #2: #3 \rarrow #4}

\newcommand{\ctxequiv}{\simeq_{\mathsf{ctx}}}
\newcommand{\Values}{\mathsf{Values}}

% DenSem
% \newcommand{\sem}[1]{\llbracket #1 \rrbracket}
\newcommand{\subfin}{\subseteq_{\fin}}
\newcommand{\inj}{\mathbf{inj}}
\newcommand{\detsem}[1]{[#1]}
\newcommand{\List}{\mathrm{List}}
\newcommand{\sIntervention}{\mathrm{Intv}}
\newcommand{\sP}[2]{\mathbf{P} \ #1 \ #2}
\newcommand{\sPP}{\mathbf{P}}
\newcommand{\sVal}{\mathbf{Val}}
\newcommand{\gradecatsymb}{\mathbb{G}}
\newcommand{\gradecat}[1]{\gradecatsymb_{#1}}
\newcommand{\sdist}{\mathrm{dist}}
\newcommand{\sintervene}{\mathrm{intervene}}

\DeclareMathOperator*{\dom}{dom}

% New language
\newcommand{\sCaus}{\mathrm{Caus}}
\newcommand{\Caus}[2]{\sCaus \ #1 \ #2}
% \newcommand{\Maybe}{\mathsf{Maybe}}
\newcommand{\opnew}{\mathsf{new}}
\newcommand{\Caustic}{\mathsf{\lambda_C}}
\newcommand{\cdtyping}[3]{#1 \vdash^d_{\Caustic} #2: #3}
\newcommand{\crtyping}[4]{#1 \vdash^r_{\Caustic} #2: #3 \rarrow #4}
\newcommand{\cptyping}[4]{#1 \vdash^{p}_{\Caustic} #2 :#3 \& #4}
\newcommand{\cptypingreordered}[4]{\cptyping{#2}{#3}{#4}{#1}}

% Compilation and Logical Relation
\DeclareMathOperator*{\len}{dom}

\newcommand{\Maybe}[1]{\mathrm{Maybe}({#1})}
\newcommand{\sem}[1]{\llbracket{#1}\rrbracket}
\newcommand{\Just}[1]{\mathrm{Just}({#1})}
\newcommand{\Nothing}{\mathrm{Nothing}}

% Dario's Macros
\usepackage{bm}
\renewcommand{\vec}{\bm}

\newcommand{\todo}[1]{\textcolor{blue}{[TODO: #1]}}
\newcommand{\prompt}[1]{\textcolor{purple}{[PROMPT: #1]}}

\newcommand{\cat}[1]{\mathbb{#1}}
\newcommand{\catname}[1]{\mathbf{#1}}

\newcommand{\set}{\mathbf{set}}
\newcommand{\Set}{\catname{Set}}

\newcommand{\concat}{+\!\!+}

\newcommand{\defeq}{\coloneq}


\newcommand{\C}{\cat C}
\newcommand{\G}{\mathcal G}

\newcommand{\op}{\mathrm{op}}


\newcommand{\thought}[3]{\marginpar{\tiny\itshape\raggedright\textcolor{#2}{\textbf{#1:} #3}}\ignorespaces} \marginparwidth=1.4cm \marginparsep=0cm
\newcommand{\jeremy}[1]{\thought{Jeremy}{green!20!black}{#1}}

\newcommand{\IntvInstr}{\mathsf{IntvInstr}}
\newcommand{\oplookup}{\mathsf{lookup}}
\newcommand{\oprestrict}{\mathsf{restrict}}
\newcommand{\FinSet}{\catname{FinSet}}
\newcommand{\Bij}{\catname{Bij}}
\newcommand{\Type}{\mathbf{Type}}

\newcommand{\vpure}{\mathbf{pure}}
\newcommand{\vapp}{\mathbin{\langle * \rangle}}
% \DeclareMathOperator*{\vapp}{\langle * \rangle}
\newcommand{\fmap}{\mathbin{\langle \$ \rangle}}

\newcommand{\colorone}[1]{\textcolor{purple}{#1}}
\newcommand{\colortwo}[1]{\textcolor{blue}{#1}}
\newcommand{\ParIndex}{\mathsf{ParIndex}}
\newcommand{\Env}{\mathbf{Env}}

\newcommand{\inl}{\mathbf{inl}}
\newcommand{\inr}{\mathbf{inr}}

\newcommand{\note}{\textcolor{purple}}
\newcommand{\cur}{\mathbf{cur}}
\newcommand{\uncur}{\mathbf{uncur}}

\newcommand{\intvPt}{\mathsf{intvPt}}
\newcommand{\Endo}{\mathcal{E}}
\newcommand{\sEndo}{\mathsf{Endo}}

\newcommand{\arrow}{\rightsquigarrow}
\newcommand{\arrarr}{\mathsf{arr}}
\newcommand{\arrcomp}{\ggg}
\newcommand{\arrfirst}{\mathsf{first}}


\newcommand{\DarioMonad}{\mathsf{T}_{\mathsf{CBN}}}
\newcommand{\oprun}{\mathsf{run}}
\newcommand{\dolet}{\mathbf{let}}
\newcommand{\doletexpr}[3]{\dolet \ #1 \leftarrow #2 \ \doin \ #3}
\newcommand{\doletexprParam}[4]{\dolet_{#1} \ #2 \leftarrow #3 \ \doin \ #4}
\newcommand{\doin}{\mathbf{in}}

\newcommand{\typing}[3]{#1 \vdash #2: #3}
\newcommand{\wtyping}[4]{#1 \vdash_{#2} #3: #4}
\newcommand{\polytyping}[3]{#1 \vdash_{[-]} #2: #3}
\newcommand{\TMon}[2]{T_{#1}\left(#2 \right)}

\newcommand{\catIntv}{\mathbf{IntvStruct}}

\newcommand{\bOne}{\mathds{1}}
\newcommand{\newEndo}{\mathsf{endo}}

\newcommand{\Meas}{\mathsf{Meas}}
\newcommand{\perWorld}{\mathsf{perWorld}}
\newcommand{\valueAt}{\mathsf{valueAt}}

\newcommand{\PSh}{\mathbf{PSh}}
\newcommand{\PShBij}{\PSh(\Bij)}

\newcommand{\LangName}{\mathsf{\lambda_C}}
\newcommand{\World}{\mathsf{World}}

\newcommand{\dwtyping}[4]{#1 \mid #2 \vdash^d #3: #4}
\newcommand{\pwtyping}[5]{#1 \mid #2 \mid #3 \vdash^p #4: #5}
% \newcommand{\rtyping}[4]{#1 \vdash #2: #3 \rarrow #4}

\newcommand{\Any}{\mathsf{Any}}
\newcommand{\intvConstr}{\mathsf{intvConstr}}
\newcommand{\twin}{\mathsf{twin}}

\newcommand{\Functions}{\mathsf{Functions}}
\newcommand{\Distributions}{\mathsf{Distributions}}
\newcommand{\Constants}{\mathsf{Constants}}

\newcommand{\rone}[1]{{#1}^{(1)}}
\newcommand{\rtwo}[1]{{#1}^{(2)}}
\newcommand{\rthree}[1]{{#1}^{(3)}}

\newcommand{\Lone}{L_1}
\newcommand{\Ltwo}{L_2}
\newcommand{\Lthree}{L_3}

\newcommand{\Sbs}{\catname{Sbs}}
\newcommand{\Qbs}{\catname{Qbs}}

\newcommand{\typingone}[4]{#1 \vdash_{1}^{#4} #2: #3}
\newcommand{\etypingone}[3]{\typingone{#1}{#2}{#3}{\mathbf{e}}}
\newcommand{\ttypingone}[3]{\typingone{#1}{#2}{#3}{\mathbf{t}}}
\newcommand{\qtypingone}[3]{\typingone{#1}{#2}{#3}{\mathbf{q}}}

\newcommand{\typingtwo}[5]{#1 \mid #2 \vdash_{2}^{#5} #3: #4}
\newcommand{\etypingtwo}[4]{\typingtwo{#1}{#2}{#3}{#4}{\mathbf{e}}}
\newcommand{\ttypingtwo}[4]{\typingtwo{#1}{#2}{#3}{#4}{\mathbf{t}}}
% \newcommand{\qtypingtwo}[4]{\typingtwo{#1}{#2}{#3}{#4}{q}}

\newcommand{\etypingthree}[4]{#1 \mid #2 \vdash_{3}^{\mathbf{e}} #3: #4}
\newcommand{\ttypingthree}[6]{#1; #2 \mid #3; #4 \vdash_{3}^{\mathbf{t}} #5: #6}
% \newcommand{\qtypingtwo}[4]{\typingtwo{#1}{#2}{#3}{#4}{q}}
\newcommand{\Gammaexo}{\Gamma_{\mathsf{exo}}}
\newcommand{\Gammaendo}{\Gamma_{\mathsf{endo}}}
\newcommand{\Deltaexo}{\Delta_{\mathsf{exo}}}
\newcommand{\Deltaendo}{\Delta_{\mathsf{endo}}}

\newcommand{\typingHO}[3]{#1 \vdash_{\mathsf{HO}} #2: #3}


\newcommand{\scoringquery}[4]{\letbind{#1}{#3}{\opscore(#2)};\ \return{#4}}

\newcommand{\xletbind}[3]{\mathbf{let \ exo} \ #1 = #2 \ \mathbf{in} \ #3}
\newcommand{\nletbind}[3]{\mathbf{let \ endo} \ #1 = #2 \ \mathbf{in} \ #3}


\newcommand{\opdefault}{\mathbf{default}}
\newcommand{\opcondition}{\mathsf{condition}}
\newcommand{\opapplyIntv}{\mathsf{applyIntv}}
\newcommand{\forget}{\mathbf{forget}}
\newcommand{\semSCM}[1]{\sem{#1}_{\mathsf{SCM}}}

% Section 4
\newcommand{\IntvConstrName}{\mathbf{IntvC}}
\newcommand{\semone}[1]{\sem{#1}_{\Lone}}
\newcommand{\semtwo}[1]{\sem{#1}_{\Ltwo}}
\newcommand{\semthree}[1]{\sem{#1}_{\Lthree}}
\newcommand{\semho}[1]{\sem{#1}_{\LHO}}
\newcommand{\liftedEndo}{\widehat{\Endo}}
\newcommand{\sliftedEndo}{\sEndo}
\newcommand{\lift}{\mathsf{lift}}
\newcommand{\liftEndo}{\mathsf{endo}}
\newcommand{\catp}{\mathbf{P}}
\newcommand{\App}{\mathsf{App}}
\newcommand{\intvLabel}{\mathsf{intvLabel}}
\newcommand{\ulBij}{\underline{\Bij}}

% Section 5
\newcommand{\LHO}{L_{\mathsf{HO}}}
\newcommand{\LthreeHOcompile}[1]{\llparenthesis #1 \rrparenthesis}
\newcommand{\IntvMonadSymb}{\mathcal{I}}
\newcommand{\IntvMonad}[2]{\IntvMonadSymb_{#1}\left(#2\right)}
\newcommand{\Twinned}{\mathcal{T}}
% \newcommand{\TwinnedIntvT}[1]{\Twinned_{#1} \IntvMonadSymb}
\newcommand{\sTwinned}{\mathsf{Twinned}}
\newcommand{\sIntvMonadSymb}{\mathsf{IntvT}}
\newcommand{\sIntvMonad}[2]{\sIntvMonadSymb_{#1}\left(#2\right)}
\newcommand{\opnewEndo}{\mathsf{newEndo}}
\newcommand{\unzip}{\mathsf{unzip}}
\newcommand{\zip}{\mathsf{zip}}

% commenting
\colorlet{jyColour}{blue}

\makeatletter
\patchcmd{\@@addmarginpar}{\ifodd\c@@page}{\ifodd\c@@page\@@tempcnta\m@@ne}{}{}
\makeatother
\reversemarginpar
\newcommand{\sidenote}[3]{\marginpar{\tiny\itshape\raggedright\textcolor{#2}{\textbf{#1:} #3}}\ignorespaces} \marginparwidth=1.4cm \marginparsep=0cm
\newcommand{\jy}[1]{\sidenote{jy}{jyColour}{#1}}


% SCMs
\DeclarePairedDelimiterX{\parensbar}[1]
{(\!\delimsize\vert}
{\delimsize\vert\!)}
{#1}
\newcommand{\compileSCM}[1]{\ensuremath{\parensbar{#1}}}
\newcommand{\scmIntervene}[2]{\ensuremath \mathsf{do}_{#1, #2}}
\newcommand{\scmTwin}[1]{\ensuremath \mathsf{twin}_{#1}}
\newcommand{\IntClass}{\ensuremath{\mathsf{Intv}}}
\newcommand{\DefaultClass}{\ensuremath{\mathsf{Default}}}

\newcommand{\commentcolor}{\textcolor{codegreen}}
\newcommand{\highlightcolor}{\textcolor{blue}}
\newcommand{\SCM}{\mathcal{M}}
